\chapter{Limitations of DCF}

No framework is universal. DCF has boundaries, failure modes, and contexts where it may not apply. Intellectual honesty requires acknowledging these limitations.

\section{Scope Limitations}

\textbf{DCF is designed for knowledge work.} It assumes:
\begin{itemize}
    \item Tasks that benefit from iterative refinement
    \item Users capable of metacognitive reflection
    \item Time for dialogue (not emergency response)
    \item Access to capable LLM systems
\end{itemize}

\textbf{DCF may not suit:}
\begin{itemize}
    \item Real-time decision systems requiring instant responses
    \item Highly structured tasks with deterministic solutions
    \item Users without foundational domain knowledge
    \item Contexts where AI interaction is inappropriate
\end{itemize}

\section{Cognitive Load Concerns}

DCF requires significant cognitive engagement. This creates risks:

\begin{itemize}
    \item \textbf{Decision fatigue}: Constant Socratic questioning can exhaust practitioners
    \item \textbf{Analysis paralysis}: Over-refinement may prevent action
    \item \textbf{Diminishing returns}: Not every interaction warrants full DCF engagement
    \item \textbf{Expertise threshold}: Beginners may lack the knowledge to evaluate AI outputs effectively
\end{itemize}

\textbf{Mitigation}: Apply DCF selectively. Reserve deep engagement for high-stakes decisions; trust automation for routine tasks.

\section{Dependency Risks}

\begin{itemize}
    \item \textbf{Skill atrophy}: Over-reliance on AI scaffolding may weaken independent capabilities
    \item \textbf{Confirmation bias}: The Thinking Mirror can reflect back what you want to see
    \item \textbf{Epistemic outsourcing}: Delegating too much judgment to AI
    \item \textbf{Context collapse}: Important nuances lost in human-AI translation
\end{itemize}

\section{Technical Limitations}

DCF inherits the limitations of underlying LLM technology:

\begin{itemize}
    \item \textbf{Hallucination}: AI may generate plausible but false information
    \item \textbf{Knowledge cutoffs}: Training data has temporal boundaries
    \item \textbf{Context window limits}: Long conversations lose early context
    \item \textbf{Inconsistency}: Same prompt may yield different results
    \item \textbf{Opacity}: Extended thinking is not fully transparent
\end{itemize}

\section{Measurement Challenges}

DCF effectiveness is difficult to quantify:

\begin{itemize}
    \item ``Clarity gained'' is subjective
    \item Collaboration quality lacks standard metrics
    \item Long-term cognitive development hard to attribute
    \item Comparison to counterfactuals (without DCF) is speculative
\end{itemize}

\section{When DCF Fails}

Observed failure patterns:

\begin{table}[H]
\centering
\begin{tabular}{ll}
\toprule
\textbf{Failure Mode} & \textbf{Symptom} \\
\midrule
Socratic theater & Going through motions without genuine inquiry \\
Infinite refinement & Never reaching ``good enough'' \\
Mirror narcissism & Using AI to confirm existing beliefs \\
Complexity addiction & Over-engineering simple problems \\
Abstraction drift & Losing connection to concrete outcomes \\
\bottomrule
\end{tabular}
\caption{DCF Failure Modes}
\end{table}

